% The overview: what ratpoints computes, the four parts of the development, the specification and
% the main theorems in words, how the proofs go, what one has to trust, what is not verified, and
% how to read the blueprint.
\chapter{Overview}\label{chap:overview}

ratpoints is a C program, built around a library function, that finds the rational points of
bounded height on the curves $y^2 = F(x, z)$, where~$F$ is a binary form of even degree with integer
coefficients, by sieving. This blueprint describes a Lean~4 development that verifies it: an
executable model of ratpoints in Lean, which follows the C code component by component; proofs
that the model meets a short specification; and a link from the model to the C code. This chapter
states the specification and the main theorems in words, each statement linked to the node where
it is stated precisely, with its Lean names; it says how the proofs go, what one has to trust and
what is not verified, and how the rest of the blueprint is organized.

\section{What ratpoints computes}\label{sec:overviewRatpoints}

The input is a list $c = [c_0, c_1, \dots, c_n]$ of integers, the coefficients of the polynomial
$f = \sum_i c_i x^i$, constant term first, and a height bound~$H$. Let~$d$ be the degree of~$f$,
the largest index with $c_d \ne 0$, and $D = d + (d \bmod 2)$ the even number at or just above it.
The curve is $y^2 = F_c(x, z)$ with the binary form
\[ F_c(x, z) = \sum_{i=0}^{d} c_i\, x^i z^{D - i} \]
of even degree~$D$. Its points lie in the weighted projective plane with the weights
$(1, D/2, 1)$, where $(x : y : z)$ and $(-x : (-1)^{D/2} y : -z)$ are the same point. ratpoints
lists the points $(a : y : b)$ with integer coordinates, normalized so that $\gcd(a, b) = 1$ and
$b > 0$, or $(a, b) = (1, 0)$ for a point at infinity, whose height $\max(|a|, b)$ is at most~$H$;
a point with $y \ne 0$ comes with both signs of~$y$. We call~$a$ the numerator and~$b$ the
denominator of the point, after the rational number~$a/b$. For instance,
\verb|ratpoints '-2 0 0 1' 130 -q| lists the points of height at most~$130$ on $y^2 = x^3 - 2$,
among them $(129 : 3830 : 100)$, as $3830^2 = 100\,(129^3 - 2 \cdot 100^3)$.

The search is done by the library function \verb|find_points|. It takes its arguments in a
structure \verb|ratpoints_args| (\csource{ratpoints.h}{362}{400}): the coefficients and the height
bound; flags that switch off the exact check of the points, the second point $(a : -y : b)$ of each
pair, the reversal of the polynomial and the Jacobi test; a restriction of the search to intervals
for~$x/z$ and to a range of the denominators; the depth of a Sturm step; and parameters that tune
the sieve. It returns no list of points. It hands each point that it finds to a function of its
caller, the \emph{callback}, together with the caller's data, adds up the values that the calls
return, and stops as soon as a call asks it to (\csource{ratpoints.h}{436}{437}). The program
\verb|ratpoints| reads the coefficients, the height bound and its options from the command line,
calls the library with a callback that prints the points in a format chosen by the options, and
prints a report around them.

Between the input and the points the C code does a good deal of work, and the chapters on the model
follow it. It checks the input and refuses some binary forms that are not squarefree. It may
replace the polynomial by its reversal, which exchanges~$x$ and~$z$, where it expects the reversed
curve to be searched faster. A Sturm step, unless it is switched off, refuses the other binary
forms that are not squarefree and narrows the search intervals to those on which~$f$ may take a
nonnegative value. It decides which denominators can carry a point: by the 2-adic information, by
the prime factors of the leading coefficient when~$d$ is odd, and by forbidden divisors and a
Jacobi symbol when~$d$ is even. It ranks small primes and some of their powers and products by a
cost model computed in doubles and chooses the moduli of the three stages of its sieve; for each
modulus and residue, a row of bits says which numerators can give a square value of the form
modulo the modulus. For each denominator it and-s these rows into bit arrays of numerators, packed
by a stride that the 2-adic information allows, with vector instructions; it tests the survivors
for common factors with the denominator and against further primes, and checks the rest exactly.
The loop over the denominators runs in blocks, between which ratpoints may correct its choice of the
moduli by what the sieve has counted, and which it may hand to a pool of threads.

The C code is that of ratpoints~3.1.1. Every line number in this blueprint refers to one commit of
its repository, the one that the record of the correspondence pins (Chapter~\ref{chap:ccode}).

\section{The four parts}\label{sec:overviewParts}

The development has four parts, in separate folders. The specification and the model import
neither each other nor the proofs, and the proofs import both; \verb|scripts/check_layout.py|
checks the imports.
\begin{enumerate}
  \item The \emph{specification} (\verb|Ratpoints/Spec|) states what ratpoints computes: the
    binary form, the normalized points of height at most~$H$, the search region that the caller
    may give, Sturm's theorem as a proposition, and the list of the points in the order in which
    ratpoints finds them. It consists of four short files that use Mathlib and nothing of the
    model, so that they can be read on their own (Chapter~\ref{chap:spec}).
  \item The \emph{model} (\verb|Ratpoints/Model|) is an executable version of ratpoints in Lean.
    It mirrors the hardware-independent structure of the C code function by function: the same
    loops, tables, bit arrays of 64-bit words and index arithmetic. It is plain code on Lean core
    and \verb|Std|, without Mathlib and without contracts (its only proofs are a few decided facts
    about its constants, the termination of one recursion and the bounds of some array accesses),
    and forms a library of its own; the program \verb|ratpoints-lean| imports it alone and has the
    command line of ratpoints (Chapters~\ref{chap:setup} to~\ref{chap:output}).
  \item The \emph{proofs} (\verb|Ratpoints/Proof|) show that the model meets the specification.
    Each component of the model has a contract, and the main theorems put the contracts together
    (Chapter~\ref{chap:theorems}). Statements of general interest that Mathlib lacks, on
    homogenizations, Sturm sequences, sign changes, valuations and the like, are stated at
    Mathlib's generality in \verb|Ratpoints/Mathlib|.
  \item The \emph{link to the C code} (Chapter~\ref{chap:ccode}): a record of which Lean
    declaration models which lines of C, kept in paragraphs \verb|C:| of the docstrings of the
    model and, for the pin, the excluded lines and the lock, in the folder \verb|correspondence|,
    checked against the C code at the pinned commit; and differential tests that run
    \verb|ratpoints-lean| and the C program on the same inputs and compare what they print.
\end{enumerate}
The model does without Mathlib so that the executable and the differential tests depend neither on
Mathlib nor on the proofs. The specification does without the model so that one can read what the
theorems claim without reading how ratpoints works.

\section{The specification}\label{sec:overviewSpec}

The specification is the one part whose correctness no proof can check. Wherever there was a
choice, it therefore copies what the C code computes, and Chapter~\ref{chap:spec} compares each of
its clauses with ratpoints~3.1.1.

\emph{Points.} The binary form~$F_c$ above is Definition~\ref{def:binaryForm}. A triple
$(x, y, z)$ of integers is a \emph{normalized point of height at most~$H$} on the curve
(Definition~\ref{def:IsPoint}) when
\[ y^2 = F_c(x, z), \qquad \gcd(x, z) = 1, \qquad z > 0 \text{ or } (x, z) = (1, 0), \qquad
   |x| \le H, \qquad z \le H\,. \]
The last four conditions say that the pair $(x, z)$ is normalized and of height at most~$H$, which
is what the statement on the output without the exact check concludes. Remark~\ref{rem:pointsC}
compares each clause with the C code; for instance, $\gcd(0, z) = z$, so that $x = 0$ occurs only
with $z = 1$, as in the test of ratpoints for common factors.

\emph{Squarefreeness.} The refusals speak about a second binary form, the polynomial made
homogeneous of the even degree $N = 2 \lceil n/2 \rceil$ that the length of the coefficient list
allows,
\[ \Phi_c = \sum_{i=0}^{n} c_i X^i Z^{N - i} \in \QQ[X, Z] \]
(Definition~\ref{def:binaryFormPoly}), and about its squarefreeness in the sense of Mathlib. This is
what ratpoints means by a polynomial that is ``not squarefree as a binary form of even degree'':
$\Phi_c$ is squarefree exactly when~$f$ is squarefree and the leading zeros of~$c$ do not lower
$\lceil n/2 \rceil$, the two conditions that the C code tests
(Lemma~\ref{lem:squarefree_binaryFormPoly}). On every input that ratpoints accepts, $\Phi_c$
and~$F_c$ take the same values.

\emph{The search region.} The caller may restrict the search to \emph{search intervals} for~$x/z$,
closed intervals whose ends are C doubles, and to a range of the denominators from
$b_{\mathrm{low}}$ to $b_{\mathrm{high}}$. The specification reads a double by its~$64$ bits: from
them it decides whether the double is a NaN and gives its value, a rational number or~$\pm\infty$.
A list of intervals is well formed when no end is a NaN, the lower end of each interval is at most
its upper end, and each interval lies strictly below the next; this is the test by which
\verb|find_points| refuses bad arguments. A point lies in the search region when it lies at
infinity, or when its denominator lies in the range and~$x/z$ lies in one of the intervals, if
there are any (Definition~\ref{def:regionSpec}). As the C code does no arithmetic on the ends that
could decide a point, the specification can speak about their exact values
(Remark~\ref{rem:regionC}).

\emph{Sturm's theorem.} The Sturm step of ratpoints drops an interval when the numbers of sign
changes of a remainder chain agree at its ends and~$f$ is negative just inside it. That no point is
lost rests on Sturm's theorem, which Mathlib does not have, although it has the Sturm sequence and
the sign changes of a list. The specification states the theorem as a proposition, for every
squarefree real polynomial and every interval $(a, b]$ (Definition~\ref{def:sturmTheorem}), and the
statements that need it take it as a hypothesis, and only when the Sturm step is on.

\emph{The order.} The list $\mathrm{orderedPoints}$ (Definition~\ref{def:orderedPoints}) gives the
points in the order in which ratpoints finds them: for the polynomial that ratpoints searches, the
reversal of~$c$ in three cases and~$c$ itself otherwise, its points at infinity, then its
denominators in the order of the loop, the numerators of each in increasing order, $y > 0$ before
$y < 0$, each point mapped back to the caller's curve after a reversal. The denominators increase
unless the polynomial searched has odd degree and a leading coefficient~$\ell$ with $|\ell| \ne 1$
that ratpoints can factor by trial division and that has fewer than $\mathrm{maxDivisors} = 512$
squarefree divisors up to the largest denominator; then they are the numbers~$e k^2$, divisor~$e$
by divisor (Definition~\ref{def:maxDegree} and Remark~\ref{rem:orderC}).

\emph{The interface of a search.} Of the arguments in \verb|ratpoints_args|, those that change the
refusals, the order or what is delivered are the \emph{options} of a search
(Definition~\ref{def:Options}): the switch of the reversal, the depth of the Sturm step (negative
for none), the search region, \verb|RATPOINTS_NO_Y|, under which of each pair only the point with
$y \ge 0$ is delivered, and \verb|RATPOINTS_NO_CHECK|, under which the numerators that survive the
three stages of the sieve and the test for common factors, and the candidate at infinity when there
is a point at infinity, are delivered with $y = 0$ without the exact check. The \emph{options of the
sieve} only steer the choice of the moduli. A search refuses, for one of five reasons that mirror
the errors 3, 5, 4, 8 and~9 of the program, or returns its \emph{outcome}, what the program reads
afterwards.

\emph{The two references.} The reference of the delivery (Definition~\ref{def:deliver}) is a fold
of a few lines that calls the callback on the points of a list in turn, adding up the values that
the calls return, until a call asks to stop; it models no C code and is the yardstick of the
protocol theorem. The brute-force reference (Definition~\ref{def:bruteForce}) tests every candidate
of the box of height~$H$ and lists exactly the normalized points of height at most~$H$, each once
(Theorem~\ref{thm:bruteForce}); the proofs compare the search with it denominator by denominator.
What one has to read to know what the main theorems say is listed in Section~\ref{sec:comparator}.

\section{The main theorems}\label{sec:overviewTheorems}

The model has three entry points, each defined from the one before
(Definition~\ref{def:findPoints}): the search with a callback and a print channel for the report of
the library, \lname{Ratpoints.findPointsVerbose}, which the program runs; the search with a
callback, \lname{Ratpoints.findPointsWith}, which is the library function as it appears to a caller
who does not ask for the report; and $\mathrm{findPoints}(c, H, s, o)$, which returns the list of
the points or a refusal, the search with the callback that appends each point to a list and never
stops the search. The theorems hold for every coefficient list~$c$, every height bound~$H$, every
set~$s$ of options of the sieve and every set~$o$ of options of a search; they assume nothing about
the options of the sieve, which the search normalizes before it chooses its moduli. Where Sturm's
theorem is needed, it enters as the hypothesis
\[ 0 \le o.\mathrm{sturm} \;\to\; \mathrm{SturmTheorem}\,, \]
which is void when the Sturm step is off.

\emph{Membership} (Theorem~\ref{thm:membership}). Assume Sturm's theorem when the Sturm step is on,
and let the exact check be on. If $\mathrm{findPoints}(c, H, s, o)$ returns a list~$L$, then a
triple lies in~$L$ exactly when it is a normalized point of height at most~$H$, with $y \ge 0$
under \verb|RATPOINTS_NO_Y|, that lies in the search region of~$o$. Indeed~$L$ is a permutation of
the brute-force reference filtered by these two conditions.

\emph{No duplicates} (Theorem~\ref{thm:nodup}). Every list that $\mathrm{findPoints}$ returns has
no duplicates. This needs no hypothesis at all, neither Sturm's theorem nor the exact check.

\emph{The order} (Theorem~\ref{thm:order}). Under the hypotheses of membership, $L$ is the list
$\mathrm{orderedPoints}$ of the specification, with the reversal allowed when~$o$ allows it and the
search region is the whole box of height~$H$, and with the largest denominator of the search as the
bound of the denominators, filtered by \verb|RATPOINTS_NO_Y| and the search region. As this list
does not depend on~$s$, the options of the sieve do not change the output with the exact check on;
nor does the depth of the Sturm step, which the theorem also states in the form that the search
with the step switched off returns the same list (this part is proved, not certified).

\emph{The refusals} (Theorem~\ref{thm:refusals}). With the Sturm step on, $\mathrm{findPoints}$
refuses exactly when~$c$ has fewer than two or more than $\mathrm{maxDegree} + 1 = 101$ entries,
$H = 0$ or $H \ge 2^{63}$ (not a C \verb|long|), $\Phi_c$ is not squarefree, or the search
intervals are not well formed. With the Sturm step off, the condition on~$\Phi_c$ is replaced by the
two tests that ratpoints makes without the remainder chain of the Sturm step, on the leading zeros
and in the reversal check. A search that refused every input would contradict these statements, so
that membership and the order are not vacuous.

\emph{The protocol} (Theorem~\ref{thm:protocol}). For every callback and its data, the search with
the callback refuses exactly when $\mathrm{findPoints}$ refuses, and otherwise gives the data and
the count of the delivery of the list of $\mathrm{findPoints}$ to the callback by the reference of
the delivery. So the calls of the callback are those for the points of that list, in their order,
up to and including the first call that stops the search, and the statements on the list carry over
to every callback, the callback of the program among them.

\emph{Without the exact check} (Theorem~\ref{thm:noCheck}). Under \verb|RATPOINTS_NO_CHECK| every
element of the list is a triple $(x, 0, z)$ whose pair $(x, z)$ is normalized, of height at
most~$H$ and in the search region; assuming Sturm's theorem when the Sturm step is on, every
normalized point $(x, y, z)$ of height at most~$H$ in the search region gives the element
$(x, 0, z)$; and the list that the search finds with the exact check on and \verb|RATPOINTS_NO_Y|
set, with~$0$ in place of each~$y$, is a sublist, so that the points come in the same order. Which
other pairs pass the sieve depends on the moduli and is not specified.

\emph{What the program prints} (Theorem~\ref{thm:cliRun}). For a command line without \verb|-v|
that it reads, \verb|ratpoints-lean| prints the head of its report, the points of the list of
$\mathrm{findPoints}$ that its callback keeps (under \verb|-1| the first), in their order and in the
format of the command line, between the strings that the command line sets, and the end of the
report (under \verb|-q| only the points and the strings). Without \verb|-x|, and assuming Sturm's
theorem when the Sturm step is on, the points kept are exactly the normalized points of height at
most~$H$ in the search region of the command line, with the exact values of the ends that it gives,
with $y \ge 0$ under \verb|-y| and none at infinity under \verb|-i|.

\emph{The blocks, the lag and the threads} (Theorems~\ref{thm:blocks}, \ref{thm:lag}
and~\ref{thm:threads}). The run delivers the points of its blocks in order, each block sifted with
moduli that depend neither on the callback nor on the options of the output; a block reads only
corrections of the moduli made at least $\mathrm{adaptLag} + 1$ blocks before it; and blocks sifted
apart and delivered in their order give what the run delivers. So the output of the model, which
sifts the blocks in turn on one thread, is that of ratpoints with any number of threads, provided
that its pool of threads behaves as its code says (Section~\ref{sec:theoremsBlocks}).

Comparator certifies eleven statements: the two on the brute-force reference, membership, no
duplicates, the two refusal statements, the order, the three statements without the exact check, and
the protocol (Section~\ref{sec:comparator}). The others are proved in the same library, but one has
to read their Lean statements to know what they say: the forms of membership with the permutation
and the finite set of the points, the part of the order theorem on the depth of the Sturm step, the
theorems on the program, and the block, lag and thread lemmas.

\section{How the proofs go}\label{sec:overviewProofs}

All the statements on the points come from one description of what a search delivers
(Section~\ref{sec:theoremsProofs}). A test of the numerators over the denominators is \emph{sound}
when it lets through the numerator of every point. The search delivers the points at infinity and
then, over the denominators of its loop in their order and the numerators of its search region in
increasing order, what the exact check finds for the numerators that pass a sound test
(Lemma~\ref{lem:searchList}). A sound test loses no point, and the exact check keeps only points;
so with the exact check on, the list of a search is the brute-force reference, filtered by the
region and by \verb|RATPOINTS_NO_Y|, in the order of the loop and mapped back after a reversal
(Lemmas~\ref{lem:searchList} and~\ref{lem:findPointsList}). The test is fixed before the
callback and the options of the output come into play, so that one test serves for all of them;
this gives the protocol and the comparison of the searches with and without the exact check.

The chapters on the model supply the pieces: the refusals in the order of the program, and the
reversal, which maps the points of the reversed curve bijectively onto those of the caller's curve
(Chapter~\ref{chap:setup}); the exact ranges of the numerators and the Sturm step, which keeps every
point under Sturm's theorem (Chapter~\ref{chap:region}); the soundness of every sieve entry that the
choice of the moduli can make and of every row that the run reads (Chapters~\ref{chap:moduli}
and~\ref{chap:rows}); the tests on the denominators and their order
(Chapter~\ref{chap:denominators}); the contract of the sieve of one denominator
(Definition~\ref{def:sift}); and the block lemma, which puts the denominators together
(Theorem~\ref{thm:blocks}). The two exits before the sieve, at the 2-adic information and at a
prime modulo which the curve has no points, are taken only when the curve has no points at all
(Lemmas~\ref{lem:twoAdic_point} and~\ref{lem:exitByPrime}).

Of the moduli the proofs need nothing but their soundness; so the theorems hold for every choice of
the moduli and every correction between the blocks, and the cost model, computed in doubles, enters
no proof (Chapter~\ref{chap:moduli}). The contract of each component with a loop is an
\verb|@[spec]| theorem, proved with \verb|vcgen|, the program logic of Lean core, from loop
invariants that are named definitions. The contracts of the sieve and of the rows use Lean core and
\verb|Std| only; those that reason modulo a prime, about valuations, Jacobi symbols or Sturm
sequences use Mathlib.

\section{What one has to trust}\label{sec:overviewTrust}

\emph{The specification.} No proof can show that the specification describes ratpoints; one has to
read it. Chapter~\ref{chap:spec} compares each clause with the C code, and three of the four files
of the specification check examples of their definitions by \verb|decide| against outputs of the C
program. The specification is written with definitions of Lean core, Batteries and Mathlib, whose
meaning one takes from these libraries: the greatest common divisor and the integer square root,
\verb|Squarefree| and \verb|IsSquare|, the prime factors of a number, the sublists of a list, the
extended rationals \verb|WithBotTop|, and, in Sturm's theorem, the Sturm sequence
\verb|Polynomial.sturmSeq|, the sign variations \verb|List.signVariations| and the cardinality
\verb|Set.ncard|. With the specification one reads the types of the model that the statements
mention and the reference of the delivery (Section~\ref{sec:comparator}).

\emph{The kernel and comparator.} Comparator, a checker for Lean proofs built into Lake, takes a
challenge that restates the specification verbatim and states the eleven certified theorems with
\verb|sorry|, and a solution that proves them by the theorems of the library. It accepts the
solution only when each theorem states the same as in the challenge, the full bodies of the
definitions that the statements reach included, the proofs use no axioms beyond \verb|propext|,
\verb|Quot.sound| and \verb|Classical.choice|, and, under \verb|--paranoid|, six kernels accept them
(Section~\ref{sec:comparator}). The project uses no tactic that adds an axiom, such as
\verb|native_decide| or \verb|bv_decide|.

\emph{The executable.} The statements are about the definitions of the model, and
\verb|ratpoints-lean| is these definitions compiled by the compiler of Lean, with its runtime and
the primitives of Lean core on machine words, arrays and floating-point numbers. That the executable
runs the model is not certified. What is checked is that the model contains no \verb|partial|,
\verb|unsafe|, \verb|implemented_by| or \verb|extern|, nor any experimental option, which could let
the compiled code differ from the definitions that the proofs see, and that the program imports the
model alone (Section~\ref{sec:ccodeCI}). The theorems on the program speak about its pure core
\lname{Ratpoints.Cli.run}; the function \lname{Ratpoints.Cli.main} around it writes the text once
and, where ratpoints is ended by \verb|SIGPIPE| when the reader of its output is gone, exits with
the status~$141$ that the shell reports for that (Chapter~\ref{chap:output}).

\emph{Sturm's theorem.} Membership, the order and the statement that every point appears without
the exact check assume Sturm's theorem when the Sturm step is on. It is a hypothesis, not an axiom:
these statements hold outright with the step off, and with the step on they say something because
the theorem, stated with Mathlib's Sturm sequence and sign variations, is true, the one mathematical
fact taken on trust; a proof of it, in Mathlib or here, would discharge the hypothesis without
changing a statement. No other certified statement needs it.

\emph{The cost model.} The search ranks its candidate moduli by a cost model that computes in
Lean's \verb|Float|, as the C code computes in doubles, with the \verb|log| and \verb|pow| of the
C library of the platform. No contract of the choice, of the blocks or of the correction reads the
value of a \verb|Float|: each states only what the natural numbers of the C code do, which holds
whatever the doubles are, NaN included (Remark~\ref{rem:moduliFloats}). So the theorems hold for
every choice of the moduli that the search can make, and the floating-point arithmetic of the
platform has to be trusted only for the agreement of the choice with that of the C program.

\section{What is not verified}\label{sec:overviewNotVerified}

\emph{The link to the C code is not a proof.} Lean has no semantics of C, and no theorem of this
blueprint speaks about the C code. The link rests on three things (Chapter~\ref{chap:ccode}): the
structure of the model, which mirrors the C code component by component, so that a reader can set
each definition beside its lines of C; the record, in which every line of code of the ten C files
lies in exactly one range, modelled by the declaration whose docstring names it or excluded with a
kind and a reason, each range locked with the hash of its text and its owners and checked by CI
against the pinned commit, and in which paragraphs \verb|Deviation:| say where the model departs
from the C code (Sections~\ref{sec:ccodeRecord} and~\ref{sec:ccodeDeviations}); and the
differential tests, which run both programs on the same command lines and compare the points in
their order, the output byte for byte, the reports and the exit statuses: $5397$ comparisons in the
default suites and~$13$ and~$92$ in two further suites, none of which fails
(Section~\ref{sec:ccodeTests}). As glibc's \verb|log| and \verb|pow| are not correctly rounded and
depend on the processor, the comparisons that depend on the choice of the moduli are made with both
programs on one machine.

\emph{What the model leaves out.} The record excludes, by kind, the pool of threads, the management
of memory, the parts of the interface of the library that only another caller meets, such as a
work space of the caller's own, the instrumentation for timing, counting and debugging, and lines
without meaning for the model: includes, prototypes, hints to the compiler, the frames of the
generators of tables, and macros, a field and a message that the C code never uses
(Section~\ref{sec:ccodeExclusions}); one more exclusion concerns other builds, below. The model
allocates nothing that can fail, so it has no error for lack of memory. In place of the threads it
has the block, lag and thread lemmas, whose hypotheses on the pool of threads are read off the C
code (Section~\ref{sec:theoremsBlocks}).

\emph{The other builds.} The vector code of the C code becomes words: a bit array of the default
build, a vector of~$256$ bits, is four consecutive~$64$-bit words. The builds of the C code differ
in compile-time constants, such as the width of the bit arrays and the number of bit arrays that the
first stage keeps in registers, which choose among the arms of the code. The model has these
constants as definitions with the values of the default build, and both arms where the C code has
two, but for the arm of \verb|exact_scale| for a build without a 128-bit integer type, which the
record excludes (\csource{find_points.c}{3239}{3256}). No proof depends on the values of the
constants beyond a few facts that are decided; a script checks this, outside CI, by building
copies of the whole project with eleven other sets of values (Section~\ref{sec:ccodeConstants}). So
the proofs cover the models of the builds that have \verb|unsigned __int128|, while the tests run
the default build.

\emph{The platform.} The doubles of the search region are their bits, with a function for each
operation that ratpoints applies to them, from the comparisons to \verb|frexp|, \verb|ldexp| and
\verb|nextafter|, and the program reads and prints numbers as glibc's \verb|strtod| and
\verb|printf| do (Chapters~\ref{chap:region} and~\ref{chap:output}). These functions model the
platform, not ratpoints: they are checked against gcc and glibc on the cases at the edges, and the
proofs show of them what the theorems need, but they are not proved to be what glibc computes.

\emph{What is compared but not proved.} The report of the library under \verb|-v| and the messages
of the program, the numbers of the moduli used, and the pairs that pass the sieve under
\verb|RATPOINTS_NO_CHECK| beyond what Theorem~\ref{thm:noCheck} says, are compared with the C
program by the differential tests. Nothing is proved about the work that a search does or about how
well the choice of the moduli serves it.

\section{How to read this blueprint}\label{sec:overviewReading}

Chapter~\ref{chap:spec} states the specification, and Chapter~\ref{chap:theorems} the main
theorems, with the proofs in outline and what comparator certifies; a reader who wants only to know
what is verified can stop there. Chapters~\ref{chap:setup} to~\ref{chap:output} describe the model
part by part, following a search from the input to the output: the set-up, the search region and
the Sturm step, the choice of the moduli and the blocks, the tables and the rows, the denominators,
the sieve of one denominator, and the delivery and the program. Chapter~\ref{chap:ccode} describes
the link to the C code, and Appendix~\ref{chap:correspondence} lists every recorded range of C.

The notation of this chapter is used throughout: $c = [c_0, \dots, c_n]$, $d$, $D$, $N$, $f$,
$F_c$ and~$\Phi_c$, the height bound~$H$, and $b_{\mathrm{low}}$ and $b_{\mathrm{high}}$ for the
range of the denominators; a point $(a : y : b)$ has the numerator~$a$ and the denominator~$b$.

\emph{Nodes.} The definitions, lemmas, theorems and corollaries are the nodes of the dependency
graph; each names its Lean declarations, which in the web version link to the API documentation. A
node of the model describes one component: it names the main definition and its contract, says in
words what the one computes and the other promises, precisely enough that a reader can check the
Lean statement against it, and links to the C code that the definition models. Each statement is
followed by its proof in outline. Remarks are not nodes; they compare the specification and the
model with the C code. Most helpers appear in no node; the API documentation has them all.

\emph{The dependency graph.} In the graph of the web version, a definition is a box and a lemma or
theorem an ellipse; an arrow leads from a node to each node that uses it, dashed when the statement
uses it and solid when only the proof does. Every node is green, and every statement dark green: its
proof and the proofs of all the nodes it rests on are formalized.

\emph{Links to the C code.} A link such as \csource{find_points.c}{4624}{4632} shows the
lines~$4624$ to~$4632$ of \verb|find_points.c| at the pinned commit. Inside a node, every link is
one of the ranges that the docstrings of its declarations record, as
\verb|scripts/correspondence.py| checks. The docstrings carry all the ranges and deviations of each
declaration, and Appendix~\ref{chap:correspondence} lists every range with its owners, or with the
reason for its exclusion.
